\documentclass[11pt]{article}
\usepackage[utf8]{inputenc}
\usepackage{geometry}
\usepackage{hyperref}
\usepackage{microtype}
\usepackage{charter}

\geometry{a4paper, margin=0.85in}

\hypersetup{
    colorlinks=true,
    linkcolor=blue,
    urlcolor=blue
}

\begin{document}

\noindent
\textbf{Ilias Chrysovergis} \\
Department of Electrical and Electronic Engineering \\
Imperial College London, London SW7 2AZ, UK \\
Metatopia \\
E-mail: \texttt{ilias@chrysovergis.gr} \\
Date: \today

\vspace{0.8em}

\noindent
\textbf{To:} \\
The Editor-in-Chief \\
\textit{IEEE Networking Letters}

\vspace{0.8em}

\noindent
\textbf{Re: Submission of Manuscript --- ``Adversarial Delay Amplification in Multi-Node Networks Under Timeout-Driven Retransmissions''}

\vspace{0.8em}

\noindent
Dear Editor,

\vspace{0.4em}
\noindent
I am pleased to submit the above manuscript for consideration as an original research Letter in \textit{IEEE Networking Letters}.

\vspace{0.4em}
\noindent
\textbf{Context.}
Modern tactical ad-hoc networks, autonomous swarms, and mission-critical IoT deployments operate over hostile channels subject to active jamming and packet corruption. Transport-layer retransmission protocols ensure reliability but introduce a non-linear positive feedback loop that can significantly amplify queueing delays. Despite the practical severity of this phenomenon, no existing analytical framework jointly models the coupled dynamics of adversarial disruption, timeout-governed retransmissions, and multi-hop topology.

\vspace{0.4em}
\noindent
\textbf{Key Contributions.}
\begin{enumerate}
    \setlength{\itemsep}{2pt}
    \setlength{\parskip}{0pt}
    \item \textbf{Destruction--Modification Duality:} I prove that post-service modification attacks degrade network capacity by a factor of $(1-p)^{-1}$ relative to pre-service destruction attacks, because corrupted packets consume complete server cycles.
    \item \textbf{Multi-Hop Delay Dynamics:} Using Gamma moment-matching for hypoexponential delay distributions, I characterize the super-linear delay scaling of $N$-node feedforward chains and verify that a robust stable operating envelope is preserved across operational parameters.
    \item \textbf{Mesh Resilience:} Through Erlang visit-count summations, I establish that symmetric feedback meshes distribute retransmission load uniformly across nodes, preventing single-node bottlenecking and yielding graceful delay scaling.
\end{enumerate}
All analytical models are rigorously validated against discrete-event simulations across five topological case studies, achieving relative errors below 2\% in all regimes and strictly under 0.65\% for multi-node networks.

\vspace{0.4em}
\noindent
\textbf{Suitability.}
The manuscript addresses a core problem in networking---the interaction between adversarial attacks and transport-layer retransmissions across multi-hop topologies---making it well suited for \textit{IEEE Networking Letters}. The unified framework provides actionable criteria for timeout dimensioning and capacity provisioning in contested environments.

\vspace{0.4em}
\noindent
\textbf{Declarations.}
This manuscript presents original, unpublished research and is not under consideration elsewhere. There are no conflicts of interest to declare.

\vspace{0.8em}

\noindent
Thank you for your consideration.

\vspace{0.8em}

\noindent
Sincerely,

\vspace{0.6em}
\noindent
\textbf{Ilias Chrysovergis} \\
E-mail: \texttt{ilias@chrysovergis.gr} / \texttt{iliachry@gmail.com}

\end{document}
